\documentclass[tikz,border=18pt]{standalone}
\usepackage{luatexja}
\usepackage{tikz}
\usetikzlibrary{shapes.geometric,arrows.meta,positioning,calc}

\definecolor{privgreen}{RGB}{30, 130, 60}
\definecolor{pubblue}{RGB}{20, 80, 160}
\definecolor{routerorange}{RGB}{215, 110, 20}

\begin{document}
\begin{tikzpicture}[>=Stealth, font=\small]

  % -------------------------------------------------------------
  % 左側：家庭・学校のプライベートLAN
  % -------------------------------------------------------------
  \node[draw=privgreen, fill=green!6, rounded corners=8pt, line width=1.5pt, minimum width=8.8cm, minimum height=7.2cm, align=center] (lan_box) at (-5.4, 2.0) {};
  \node[above, font=\large\bfseries, text=privgreen] at (-5.4, 4.8) {【プライベートLAN】学校・自宅のネットワーク};

  \node[draw=privgreen, fill=white, rounded corners=4pt, thick, align=center, minimum width=7.2cm, minimum height=1.3cm] (lan_pc) at (-5.4, 3.8) {
    \textbf{高校生パソコン (DHCPで自動割当)}\\[2pt]
    \footnotesize IP: \texttt{192.168.1.15} / MAC: \texttt{00:1A:2B:...}
  };

  \node[draw=privgreen, fill=white, rounded corners=4pt, thick, align=center, minimum width=7.2cm, minimum height=1.3cm] (lan_phone) at (-5.4, 2.0) {
    \textbf{高校生スマホ}\\[2pt]
    \footnotesize IP: \texttt{192.168.1.24}
  };

  \node[draw=black!70, fill=yellow!15, rounded corners=4pt, align=center, minimum width=7.2cm, minimum height=1.5cm] (lan_dns) at (-5.4, 0.0) {
    \textbf{DNSキャッシュサーバー}\\[2pt]
    \footnotesize \texttt{www.example.com} $\rightarrow$ \texttt{203.0.113.80}\\
    \scriptsize ドメイン名をIPアドレスに電話帳のように変換！
  };

  % -------------------------------------------------------------
  % 中央：ルーター（NAT / NAPT：ポート番号で振り分け）
  % -------------------------------------------------------------
  \node[draw=routerorange, fill=orange!20, rounded corners=8pt, line width=2pt, align=center, minimum width=4.0cm, minimum height=4.2cm] (router) at (0.0, 2.0) {
    \textbf{\Large ルーター}\\[4pt]
    \textbf{\large (NAT / NAPT)}\\[4pt]
    \footnotesize プライベートIP と\\
    \footnotesize グローバルIP を\\
    \footnotesize 相互変換！\\[3pt]
    \scriptsize 内外を遮断する\\
    \scriptsize 最初のファイアウォール
  };

  % -------------------------------------------------------------
  % 右側：グローバルインターネット
  % -------------------------------------------------------------
  \node[draw=pubblue, fill=blue!6, rounded corners=8pt, line width=1.5pt, minimum width=8.8cm, minimum height=7.2cm, align=center] (wan_box) at (5.4, 2.0) {};
  \node[above, font=\large\bfseries, text=pubblue] at (5.4, 4.8) {【グローバルWAN】世界中のインターネット};

  \node[draw=pubblue, fill=white, rounded corners=4pt, thick, align=center, minimum width=7.2cm, minimum height=1.5cm] (pub_ip) at (5.4, 3.6) {
    \textbf{契約プロバイダ割当のグローバルIP}\\[2pt]
    \footnotesize IP: \texttt{198.51.100.42} (世界で唯一の宛先)
  };

  \node[draw=pubblue, fill=blue!15, rounded corners=4pt, very thick, align=center, minimum width=7.2cm, minimum height=2.2cm] (web_srv) at (5.4, 0.8) {
    \textbf{\large 目的のWebサーバー}\\[3pt]
    \footnotesize IP: \texttt{203.0.113.80} (Port: 443 HTTPS)\\[2pt]
    \scriptsize パケットがルーターを経由してバケツリレーで届く！
  };

  \draw[->, very thick, privgreen] (lan_pc.east) -- (router.west |- lan_pc.east);
  \draw[->, very thick, privgreen] (lan_phone.east) -- (router.west |- lan_phone.east);
  \draw[->, very thick, pubblue] (router.east) -- (pub_ip.west);
  \draw[->, very thick, pubblue] (pub_ip.south) -- (web_srv.north);

  % 下部まとめ
  \node[draw=black!70, fill=green!20, rounded corners=6pt, line width=1.2pt, align=center, minimum width=20.0cm, inner sep=6pt] at (0, -2.4) {
    \textbf{\Large ★ ネットワークが分かるとセキュリティの「どこを守るべきか」が見えてくる！}\\[3pt]
    \small IPアドレス（住所）、MACアドレス（機器識別）、ポート番号（部屋番号）、DNS（案内所）の連携がインターネットの基礎！
  };

\end{tikzpicture}
\end{document}
